The first isomorphism is Fact \claimref{lepton.imported-duality-pair-system} as in the proof of Proposition \claimref{lepton.content-constituent-space}(3), without setting the determinants to one. For the non-isomorphism, compare $K$-weights: $\Con_{j}$ has $K$-weight $1 - j$ and $\Con_{k}^{*}$ has $-(1 - k) = k - 1$; equality needs $j + k = 2$, so $(j, k) \in \{(0,2), (1,1), (2,0)\}$, but the $SU(3)$ contents of $\Con_{0}, \Con_{1}, \Con_{2}$ are $1, 3, \bar{3}$ and those of $\Con_{2}^{*}, \Con_{1}^{*}, \Con_{0}^{*}$ are $3, \bar{3}, 1$, which never match. If $\det A \cdot \det B = 1$ on $G$ the twist is trivial and $\Con_{k}^{*} \cong \Con_{3-k}$; if not, some $(A, B) \in G$ has $\det A \cdot \det B = \mu \neq 1$, and then $\Con_{1}^{*} \cong \Con_{2} \otimes \mu^{-1}$, whose $SU(3)$ content $\bar{3}$ is matched by $\Con_{2}$ alone among the $\Con_{j}$, is not isomorphic to $\Con_{2}$ (they differ on that element), so $\Con$ is not self-dual. (The case $k = 0$ would not do: $\Con_{0}$ and $\Con_{3}$ share the content $(2,1)$, and for $G = SU(2) \times U(3)$ one has $\Con_{0}^{*} \cong \C^{2} \otimes (\det A)^{-1} = \Con_{0}$, so the dual of $\Con_{0}$ is present there although $\Con$ as a whole is not self-dual; $k = 1$ or $k = 2$, whose $SU(3)$ contents are matched by one constituent each, decide.) Since $G \supseteq SU(2) \times SU(3)$, the condition $\det A \cdot \det B = 1$ on $G$ is the condition $G \subseteq S(U(2) \times U(3))$. The script checks the $K$-weights and the dual weights, and tests the biconditional on five subgroups $G$, through the lattice of central characters trivial on $G$, including the one on which the $k = 0$ test alone would fail to decide.
